\documentclass[12pt,a4paper,oneside]{report}

% --- Basis-Pakete für Sprache und Schrift ---
\usepackage[utf8]{inputenc}
\usepackage[german]{babel}
\usepackage[T1]{fontenc}
\usepackage{helvet} 
\renewcommand{\familydefault}{\sfdefault} 

% --- Zeilenabstand & Mikrotypografie ---
\usepackage{setspace}
\setstretch{1.15} 
\usepackage{microtype} 

% --- Farben (Dezente Luftfahrt-CI) ---
\usepackage{xcolor}
\definecolor{AerospaceNavy}{RGB}{20, 40, 80} 
\definecolor{Charcoal}{RGB}{50, 50, 50}      
\definecolor{LightGrey}{RGB}{230, 235, 245}  % Für den dezenten Trennbalken

% --- Seitenränder & Layout ---
\usepackage[left=3cm,right=2.5cm,top=3cm,bottom=3cm]{geometry}
\usepackage{fancyhdr}
\pagestyle{fancy}
\fancyhf{}
\lhead{\small\textcolor{Charcoal}{\leftmark}}
\rhead{\small\textcolor{Charcoal}{Dual-Tilt-Rotor TDD --- Bundesheer SAR}}
\rfoot{\textcolor{Charcoal}{\thepage}}
\renewcommand{\headrulewidth}{0.4mm}
\renewcommand{\headrule}{\hbox to\headwidth{\color{AerospaceNavy}\leaders\hrule height \headrulewidth\hfill}}

% --- Überschriften einfärben ---
\usepackage{titlesec}
\titleformat{\chapter}[hang]{\LARGE\bfseries}{\thechapter\quad}{0pt}{\color{AerospaceNavy}}
\titleformat{\section}{\Large\bfseries}{\thesection\quad}{0pt}{\color{AerospaceNavy}}
\titleformat{\subsection}{\large\bfseries}{\thesubsection\quad}{0pt}{\color{AerospaceNavy}}

% --- Grafiken, Tabellen & Mathematik ---
\usepackage{graphicx}
\usepackage{booktabs}
\usepackage{amsmath}
\usepackage{amssymb}
\usepackage{array} 

% --- Hyperref Konfiguration ---
\usepackage{hyperref}
\hypersetup{
    colorlinks=true,
    linkcolor=AerospaceNavy, 
    filecolor=AerospaceNavy,      
    urlcolor=AerospaceNavy,  
    citecolor=AerospaceNavy,
    pdftitle={Technical Design Document - Dual-Tilt-Rotor VTOL},
    pdfauthor={DaSommsi}
}

% =======================================================
% START DES DOKUMENTS
% =======================================================
\begin{document}

\color{Charcoal} 

% --- MODERNES, LINKSBÜNDIGES DECKBLATT (OPTIMIERT) ---
\begin{titlepage}
    \flushleft 
    
    % Logo-Bereich ganz oben
    \begin{figure}[h!]
        \flushleft
        % \includegraphics[height=1.5cm]{logos/bundesheer_logo.png}
        \hspace{1cm}
        % \includegraphics[height=1.5cm]{logos/htl_logo.png}
    \end{figure}
    
    \vspace*{2.5cm}
    
    % Vertikaler Akzentbalken und Titel
    \noindent
    \begin{tabular}{!{\color{AerospaceNavy}\vrule width 4pt} p{\dimexpr\textwidth-25pt}}
        \hspace{15pt} \begin{minipage}{\linewidth}
            {\LARGE\bfseries TECHNICAL DESIGN DOCUMENT\par}
            \vspace{0.4cm}
            {\Huge\bfseries DUAL-TILT-ROTOR VTOL\par}
            \vspace{0.6cm}
            {\large\color{Charcoal}Modulares 3m-VTOL-System mit schwarmbereiter Architektur für den Katastrophenschutz, Personensuche und Echtzeit-Kartierung\par}
        \end{minipage}
    \end{tabular}

    \vfill
    \vfill
    
    % Horizontaler dünner Trennbalken vor den Metadaten
    \noindent\textcolor{LightGrey}{\rule{\textwidth}{1pt}}
    \vspace*{0.6cm}

    % Cleane, perfekt strukturierte Fußdaten (Zweispaltig mit klarer Hierarchie)
    \begin{minipage}[t]{0.5\textwidth}
        {\footnotesize\color{AerospaceNavy}\uppercase{\textbf{Autoren}}\par}
        \vspace{0.1cm}
        {\large\textbf{David Sommerer}}\par
        {\large\textbf{Sebastian Minkenberg}}\par
        
        \vspace{0.6cm} % Deutlicher Abstand zur nächsten Kategorie
        
        {\footnotesize\color{AerospaceNavy}\uppercase{\textbf{Organisation}}\par}
        \vspace{0.1cm}
        HTL-Diplomprojekt\par
    \end{minipage}
    \begin{minipage}[t]{0.45\textwidth}
        {\footnotesize\color{AerospaceNavy}\uppercase{\textbf{Kooperation \& Zertifizierung}}\par}
        \vspace{0.1cm}
        Österreichisches Bundesheer\par
        Austro Control --- SORA SAIL II/III\par
        
        \vspace{0.6cm} % Symmetrischer Abstand auf der rechten Seite
        
        {\footnotesize\color{AerospaceNavy}\uppercase{\textbf{Dokumentenstand}}\par}
        \vspace{0.1cm}
        \begin{tabular}{ll}
            Version: & 1.0.0 \\
            Datum:   & \today \\
        \end{tabular}
    \end{minipage}

    \vspace*{1cm}
\end{titlepage}

% --- Verzeichnisse ---
\tableofcontents
\newpage

% --- Abkürzungsverzeichnis ---
\chapter*{Abkürzungsverzeichnis}
\addcontentsline{toc}{chapter}{Abkürzungsverzeichnis}

\begin{tabular}{p{2.5cm} p{11cm}}
    \toprule
    \textbf{Kürzel} & \textbf{Bedeutung / Beschreibung} \\
    \midrule
    \textbf{ANSYS} & Analysis System (High-End Simulationssoftware) \\
    \textbf{CFD}   & Computational Fluid Dynamics (Strömungssimulation) \\
    \textbf{CoG}   & Center of Gravity (Schwerpunkt) \\
    \textbf{ESC}   & Electronic Speed Controller (Drehzahlregler) \\
    \textbf{FEA}   & Finite Element Analysis (Strukturanalyse) \\
    \textbf{KPI}   & Key Performance Indicator (Leistungskennzahl) \\
    \textbf{ROS 2} & Robot Operating System 2 \\
    \textbf{SAR}   & Search and Rescue (Such- und Rettungsdienst) \\
    \textbf{SORA}  & Specific Operations Risk Assessment (Risikoanalyse) \\
    \textbf{VTOL}  & Vertical Take-Off and Landing (Senkrechtstarter) \\
    \bottomrule
\end{tabular}
\newpage

% --- Automatische Einbindung der 13 Industrie-Kapitel ---
\chapter{Executive Summary \& Admin}
\label{ch:admin}
\section{Purpose \& Scope}
\section{Document Traceability (MRD)}
\section{Applicable Standards (EASA SC-Light UAS, STANAG 4703)}
\section{Configuration Management}

\chapter{Concept of Operations (CONOPS)}
\label{ch:conops}
\section{Operational Scenarios (VTOL, Transition, Fixed-Wing)}
\section{Flight Envelopes}
\section{Environmental Constraints}

\chapter{System Architecture \& High-Level Design}
\label{ch:architecture}
\section{System Element Breakdown (Air Vehicle vs. GCS)}
\section{Functional Architecture}
\section{Physical Architecture \& Integration}

\chapter{Airframe \& Aerodynamics}
\label{ch:aerodynamics}
\section{Aerodynamic Design \& Transition Physics}
\section{Structural Layout \& Materials (CFK, PA12-CF)}
\section{Mass Properties \& Center of Gravity (CoG) Shift}

\chapter{Propulsion \& Tilt-Mechanism}
\label{ch:propulsion_tilt}
\section{Powertrain Architecture}
\section{Motor Specifications \& Efficiency}
\section{Tilt Actuation System (TAS) \& Redundancy}
\section{Rotor/Propeller Design}

\chapter{Avionics \& Flight Control System (FCS)}
\label{ch:avionics_fcs}

\section{Avionics Hardware Architecture}
\label{sec:avionics_hardware_architecture}

\subsection{Flight Management Unit (FMU) \texorpdfstring{\&}{\&} Baseboard Configuration}
\label{sec:fmu_baseboard_config}

Das elektronische Herzstück der Avionik-Infrastruktur bildet eine modulare Kombination aus der Flight Management Unit (FMU) und einem robusten Trägerboard 
(v2B Standard Baseboard). Diese physische Trennung von Rechen- und Schnittstellen-Ebene ermöglicht eine hohe Wartungsfreundlichkeit sowie eine 
gezielte Optimierung der Signalwege und Vibrationsentkopplung.

\subsubsection{Processor and Sensor Redundancy}
\label{sec:processor_sensor_redundancy}

Zur Erfüllung der anspruchsvollen Ausfallsicherheitskriterien innerhalb der behördlichen SAIL-IV-Zulassung basiert 
der eingesetzte Autopilot \textit{Pixhawk 6X} (Revision 8) auf einer strikt partitionierten, hardwareseitig redundanten Architektur 
nach dem etablierten Open-Hardware-Standard \textit{FMUv6X} \cite{holybro_pixhawk6x}. Die konsequente Trennung der 
Rechen- und Sensordomänen eliminiert \textit{Single Points of Failure} (SPoF) auf Systemebene.

Die Rechenlast wird deterministisch auf zwei autarke Mikroprozessoren aufgeteilt:
\begin{itemize}
    \item \textbf{FMU-Hauptprozessor:} Ein \textit{STM32H753} 32-Bit ARM Cortex-M7 Prozessor ($480\,\text{MHz}$, $2\,\text{MB}$ Flash, $1\,\text{MB}$ RAM). 
                                       Dieser Primärprozessor verarbeitet die hochfrequenten Lageregelkreise (PID-Schleifen) des Tri-Tilt-Rotor-Mischprofils 
                                       und führt das POSIX-konforme Echtzeitbetriebssystem \textit{PX4 Autopilot} (Soll-Release $\ge 1.14.3$) aus.
    \item \textbf{IO-Koprozessor:} Ein physisch separierter \textit{STM32F103} 32-Bit ARM Cortex-M3 Prozessor ($72\,\text{MHz}$, $64\,\text{KB}$ SRAM). 
                                   Dieser überwacht im Hintergrund die Ausgänge und fungiert im Falle eines kritischen Software-Freezes des Hauptprozessors 
                                   als Hardware-Failsafe-Ebene zur zyklischen Aufrechterhaltung von Notfall-Steuerbefehlen wie dem Fallschirmauswurf.
\end{itemize}

Die Erfassung der Raumlage- und Beschleunigungsdaten erfolgt über ein dreifach redundantes Sensor-Domänen-Konzept (\textit{Triple Redundancy Domains}). 
In der Revision 8 sind drei identische \textit{ICM-45686} Trägheitsmessheiten (IMUs) integriert, welche lineare Beschleunigungen von bis zu $\pm32\,\text{g}$ 
fehlerfrei abbilden. Diese Sensoren nutzen die patentierte \textit{BalancedGyro\texttrademark}-Technologie, um temperatur- und vibrationsbedingte 
Drehratendriften inhärent zu kompensieren. 

Jede dieser drei IMU-Domänen ist auf Chipebene vollständig isoliert: Sie wird über mathematisch autarke SPI-Busse ausgelesen und von einer dedizierten, 
voneinander getrennten Spannungsregelung (\textit{Power Control}) versorgt. Ein elektrischer Kurzschluss innerhalb einer Sensorplatine führt somit nicht 
zum Gesamtausfall des Navigationssystems. Zur Erhöhung der Zuverlässigkeit bei der barometrischen Höhenbestimmung wird zudem eine diversitäre Sensorbestückung, 
bestehend aus einem \textit{ICP20100} und einem \textit{BMP388} Barometer, eingesetzt, was softwarebasierte Gleichreaktionsfehler wirksam ausschließt.

Ein kritischer Aspekt für den alpinen Einsatz bei extremen Wetterwechseln ist das integrierte thermische und mechanische Dämpfungskonzept des IMU-Boards.
 Die Sensoren befinden sich auf einer im Rumpf schwimmend gelagerten Ebene, die mittels eines custom-formulierten, hochfesten Isolationsmaterials mechanisch entkoppelt ist. 
 Dieses Material verschiebt die Resonanzfrequenz der Dämpfung gezielt in ein höheres Frequenzspektrum, wodurch hochfrequente Störvibrationen der 
 drei Brushless-Hauptmotoren vollständig gefiltert werden. Über integrierte Heizwiderstände wird das IMU-Board zudem elektronisch permanent auf einer 
 konstanten Betriebstemperatur gehalten. Dies eliminiert mathematische Sensor-Drifts beim schnellen Aufstieg in kalte Luftschichten im Hochgebirge und garantiert 
 die präzise Stützung des Extended Kalman Filters (EKF).

\subsubsection{Pinout and Port Assignment Matrix}

\subsection{Companion Computer (CC) Integration}
\label{sec:companion_computer_integration}

Die Bewältigung rechenintensiver Missionsaufgaben, die über die deterministischen Steuerungsaufgaben der FMU hinausgehen, insbesondere die Echtzeit-Bildverarbeitung der EO/IR-Nutzlast sowie die Schwarmkommunikation, 
erfolgt über einen dedizierten Companion Computer. Das System nutzt hierfür das \textit{NVIDIA Jetson Orin Nano Developer Kit} \cite{nvidia_jetson_orin_nano}. Diese Plattform kombiniert ein leistungsstarkes System-on-Module (SoM) 
mit einem vielseitigen Referenz-Trägerboard (\textit{Reference Carrier Board}). Mit einer KI-Rechenleistung von bis zu $40\,\text{TOPS}$ (bzw. bis zu $67\,\text{TOPS}$ bei optimierten generativen KI-Pipelines) 
bei einer konfigurierbaren Leistungsaufnahme von \SI{7}{\W} bis \SI{15}{\W} stellt das Modul die primäre Recheneinheit für Missionsalgorithmen innerhalb der \SI{25}{\kg}-Flugklasse dar.

\subsubsection{Carrier Board Interfaces}
\label{sec:carrier_board_interfaces}

Das integrierte \textit{Reference Carrier Board} bricht die hochpoligen Signalleitungen des SoMs auf standardisierte Schnittstellen für die Drohnenperipherie auf. 
Die physischen Anschlüsse des Companion Computers sind gemäß Hardwarespezifikation \cite{nvidia_jetson_orin_nano} wie folgt definiert:
\begin{itemize}
    \item \textbf{Gigabit Ethernet (1x GbE):} Bildet das primäre High-Speed-Netzwerk. Es dient als Kommunikationsbrücke zur FMU via \textit{micro-XRCE-DDS} im ROS~2-Netzwerk sowie zur direkten netzwerkbasierten Anbindung der EO/IR-Gimbal-Nutzlast 
                                              (Raptor / NightHawk / Seeker) zum Empfang der H.264/H.265-RTSP-Videostreams.
    \item \textbf{USB 3.2 Gen 2 (4x Type-A):} Bereitstellung von vier hochbandbreitigen Anschlüssen für sekundäre Peripheriegeräte (z.~B. externe Speichermedien oder Backup-Kamerasysteme).
    \item \textbf{MIPI CSI-2 (2x 22-pin):} Kamera-Schnittstellen für die direkte, latenzfreie Anbindung von Sekundärsensorik.
    \item \textbf{Erweiterungsanschlüsse (40-pin Header):} Stellt serielle Schnittstellen (UART, SPI, I2C, GPIO) zur Verfügung, die als physikalische Redundanzleitungen (Backup-Bridge) zur FMU geschaltet sind.
\end{itemize}

\subsubsection{Thermal Management and Cooling}
\label{sec:thermal_management_cooling}

Aufgrund der kompakten Abmessungen des Gesamtsystems (\SI{103}{\mm} $\times$ \SI{90.5}{\mm} $\times$ \SI{34.77}{\mm} inkl. Trägerboard) und einer maximalen Verlustleistung von \SI{15}{\W} erfordert die Systemintegration ein aktives thermisches Management, 
um ein thermisch bedingtes Herabsetzen der Taktfrequenz (\textit{Thermal Throttling}) bei hohen Umgebungstemperaturen zu verhindern.

Das Kühlungskonzept basiert auf einer interdisziplinären Schnittstelle zwischen Avionik-Hardware und Strukturmechanik. Das Trägerboard stellt hierfür einen dedizierten \SI{4}{-pin} Lüfteranschluss (\textit{4-pin fan header}) bereit. 
Über diesen Anschluss wird ein passgenauer, kugelgelagerter Aktivlüfter über ein PWM-Signal direkt vom Jetson-Betriebssystem in Abhängigkeit von den internen Kern-Temperaturen dynamisch geregelt. 
Die maschinenbauliche Integration des Kühlkörpers erfolgt über das strukturelle Rumpfdesign (siehe Kapitel~X), welches über aerodynamische Einlassöffnungen (\textit{NACA-Ducts}) einen kontinuierlichen Fahrtwindstrom direkt über die Aluminium-Kühlrippen leitet.


\subsection{Mission Periphery \texorpdfstring{\&}{\&} Sensor Integration}
\label{sec:mission_periphery_sensor_integration}

Die Erfassung von Umgebungs- und Positionsdaten basiert auf einer hochredundanten Sensorarchitektur. Alle missions- und navigationskritischen Peripheriemodule sind so in das Gesamtsystem integriert, 
dass ein Sensorausfall durch hard- und softwareseitige Failover-Ebenen vollständig kompensiert wird.

\subsubsection{Primary and Secondary GNSS Architecture (Dual-Sensor Setup)}
\label{sec:gnss_architecture}

Zur physischen Umsetzung der SAIL-IV-Sicherheitsanforderungen verfügt das UAS über ein duales, geometrisch und technologisch getrenntes Satellitennavigationssystem (GNSS). 
Diese Redundanz unterbindet den Totalverlust der absoluten Positionsdaten bei Hardwaredefekten oder lokalen Signalstörungen (z.~B. Abschattungen im alpinen Raum).

\begin{itemize}
    \item \textbf{Primäres GNSS (Holybro M10 GPS):} Dieses Modul basiert auf dem \textit{u-blox M10}-Chipsatz und ist über eine standardisierte 10-pin JST-GH Verbindung direkt an den \textit{GPS1}-Port 
                                                    des Pixhawk 6X Baseboards angeschlossen \cite{holybro_m10_gps}. Es dient im Normalbetrieb als primäre Quelle für die dreidimensionale Positionsbestimmung. 
                                                    Zudem beinhaltet das Gehäuse ein digitales \textit{IST8310}-Magnetometer, welches über den I2C-Bus die absolute Ausrichtung (Kompasskurs) des Luftfahrzeugs liefert.
    \item \textbf{Sekundäres GNSS (Holybro H-RTK F9P Rover 2nd GPS):} Als redundante Ebene fungiert ein hocheffizienter Multiband-Empfänger auf Basis des \textit{u-blox ZED-F9P}-Chipsatzes, welcher am \textit{GPS2}-Port der FMU angebunden ist 
                                                                      \cite{holybro_hrtk_f9p}. Durch die simultane Auswertung der Frequenzbänder L1, L2 und L5 bietet dieser Empfänger eine extrem hohe Störfestigkeit gegen Mehrwegeeffekte. 
                                                                      Bei aktiver Anbindung an eine RTK-Bodenstation über den C2-Datenlink liefert das System eine zentimetergenaue Echtzeit-Positionspräzision.
\end{itemize}

Die PX4-Autopilot-Software fusioniert die Datenströme beider GNSS-Empfänger permanent im Extended Kalman Filter (EKF) über einen dualen \textit{GPS-Blending-Algorithmus}. 
Die Gewichtung der Positionsdaten erfolgt dynamisch auf Basis der aktuellen Signalqualität (Sichtbare Satelliten, HDOP/VDOP-Werte). Weicht ein Sensor aufgrund eines Defekts oder externer Störeinflüsse mathematisch vom Erwartungshorizont ab, 
wird er vom EKF isoliert. Das System führt die Navigation ohne transiente Fluglagestörungen rein auf Basis des verbleibenden Empfängers fort.

\subsubsection{Electro-Optical / Infrared (EO/IR) Gimbal Payload}
\label{sec:eo_ir_gimbal_payload}

Die primäre Sensorik zur Durchführung der Such-, Rettungs- und Aufklärungsmissionen besteht aus einem kreiselstabilisierten, multi-spektralen EO/IR-Kamerasystem an der Bugspitze des Rumpfmonocoques. 
Um maximale Flexibilität hinsichtlich der Beschaffungslogistik, potenzieller Sponsoring-Vereinbarungen oder eines harten Budget-Limits von maximal \SI{5000}{\text{€}} zu gewährleisten, nutzt das UAS ein modulares Payload-Konzept. 
Die mechanischen Aufhängungen, die Spannungsversorgung (Wide Voltage Input) sowie die softwareseitigen Kommunikationsprotokolle sind universell ausgelegt. 

Die primäre Signalübertragung der Videostreams erfolgt netzwerkbasiert über die Ethernet-Infrastruktur direkt an das \textit{NVIDIA Jetson Orin Nano Developer Kit}. Die Videodaten werden onboard als hochkomprimierte RTSP-Streams 
(Real-Time Streaming Protocol) bereitgestellt und direkt in die gängigen KI-Inferenz-Pipelines (z.~B. YOLO-Architekturen zur Personendetektion) eingespeist. Die Steuerung erfolgt hybrid: Während die inertiale Kompensation von Flugbewegungen 
intern über das Gimbal gelöst wird, werden autonome Objektverfolgungsbefehle (\textit{Autonomous Object Tracking}) über das ROS~2-Netzwerk des Companion Computers an das System übermittelt.

Zur Realisierung dieses modularen Ansatzes sind drei Hardwarereferenzen spezifiziert:

\begin{enumerate}
    \item \textbf{High-End-Klasse (Referenzmodell: NextVision RAPTOR):} Diese Konfiguration stellt das technologische Maximum für den High-Speed-Transitionsflug dar. Bei einer Masse von $640\,\text{g}$ kombiniert das System 
                                                                        einen kontinuierlichen 80-fach-Gesamtzoom ($40\times$ optisch, $2\times$ digital) mit einem hochauflösenden, ungekühlten HD-Thermalsensor ($1280 \times 720$ Pixel) \cite{nextvision_raptor}.
                                                                         Die keilförmige Aero-Geometrie minimiert den parasitären Luftwiderstand im Fixed-Wing-Reiseflug. Dieses Modell ist primär für eine akademische Leihgabe (\textit{Loan Agreement}) 
                                                                         über das Österreichische Bundesheer vorgesehen.
    \item \textbf{Militärische Kaufklasse (Referenzmodell: NextVision NIGHTHAWK2-UZ):} Dieses System fungiert als voll operationale und gewichtsoptimierte Alternative. Mit einer Masse von lediglich $320\,\text{g}$ entlastet dieses Modul 
                                                                                       die Avionik-Massenbilanz signifikant und bietet bei kompakter zylindrischer Bauform ($90 \times 120\,\text{mm}$) einen 40-fachen Gesamtzoom sowie denselben 
                                                                                       HD-Wärmebildkanal ($1280 \times 720$ Pixel) wie das High-End-Modell \cite{nextvision_nighthawk}. Die softwareseitige API bleibt innerhalb des NextVision-Ökosystems identisch.
    \item \textbf{Wirtschaftliche Kaufklasse (Referenzmodell: Foxtech Seeker EO \& Ranging AI):} Dieses System bildet die primäre, kommerziell direkt beschaffbare Stufe innerhalb des definierten Budgetrahmens ($\approx \SI{3900}{\text{€}}$) \cite{foxtech_seeker}. 
                                                                                                 Das System nutzt ein innovatives, nicht-orthogonales 3-Achsen-Gimbal, das sich eng an die Sensoreinheit schmiegt. Dies garantiert eine geringe Stirnfläche 
                                                                                                 und exzellente Aerodynamik ohne das typische Schütteln klassischer Kugel-Gimbals bei hohen Reisegeschwindigkeiten. Die Sensorik verfügt über ein duales Sichtfeld
                                                                                                  (\textit{Dual Field of View}) mit einer Weitwinkel- und einer 30-fach-Hybridzoom-Kamera ($10\times$ optisch). Ein integrierter Laser-Entfernungsmesser 
                                                                                                  (\textit{Laser Rangefinder}) liefert zudem in Echtzeit die exakten Distanzdaten und globalen Zielkoordinaten direkt an das On-Screen-Display (OSD) 
                                                                                                  und den Jetson-Companion-Computer, unterstützt durch eine on-board AI-Mehrzielerkennung für Fußgänger und Fahrzeuge.
\end{enumerate}

\begin{table}[h!]
    \centering
    \caption{Vergleichende Trade-Off-Matrix der modularen EO/IR-Gimbal-Optionen}
    \vspace{0.2cm}
    \begin{tabular}{p{3.8cm} p{2.8cm} p{2.8cm} p{2.8cm}}
        \toprule
        \textbf{Parameter} & \textbf{NextVision RAPTOR} & \textbf{NextVision NIGHTHAWK2} & \textbf{Foxtech Seeker AI} \\
        \midrule
        \textbf{Masse} & $\SI{640}{\g}$ & $\SI{320}{\g}$ & $\approx \SI{800}{\g}$ \\
        \textbf{Aero-Geometrie} & Keilform (Exzellent) & Zylindrisch (Gut) & Nicht-orthogonal (Gut) \\
        \textbf{EO-Spezifikation} & $80\times$ Gesamt-Zoom & $40\times$ Gesamt-Zoom & 2K, $30\times$ Hybrid-Zoom \\
        \textbf{Infrarot (IR)} & $1280 \times 720$ Pixel (HD) & $1280 \times 720$ Pixel (HD) & Optional \\
        \textbf{Positionsermittlung} & Software-Raycasting & Software-Raycasting & Hardware Laser Rangefinder \\
        \textbf{Marktpreis (ca.)} & $> \SI{10000}{\text{€}}$ (Leihe) & $\approx \SI{8000}{\text{€}}$ (Kauf) & $\approx \SI{3900}{\text{€}}$ (Kauf) \\
        \textbf{SAIL-IV Eignung} & Uneingeschränkt & Uneingeschränkt & Voll operational \\
        \bottomrule
    \end{tabular}
    \label{tab:gimbal_trade_off}
\end{table}

\subsubsection{Optical Flow, LiDAR Rangefinder \texorpdfstring{\&}{\&} Vision-Based Navigation Stack}
\label{sec:navigation_failsafe_stack}

Zur Gewährleistung der primären Systemsicherheit innerhalb der behördlichen SAIL-IV-Zulassung verfügt das UAS über eine mehrstufige, redundante Navigations- 
und Driftkompensationsebene. Dieses kombinierte Sensor- und Softwaresystem dient als Ausfallschutz (\textit{Failsafe}) bei einem transienten oder permanenten 
Verlust des GNSS-Signals (z.~B. durch Signalabschattung oder gezieltes Jamming im alpinen Einsatzraum). Die Architektur ist hierarchisch nach Flughöhen gestaffelt 
und partitioniert die Rechenlast zwischen der deterministischen Flugsteuerungsebene (Pixhawk 6X) und dem übergeordneten Companion Computer 
(\textit{NVIDIA Jetson Orin Nano Developer Kit}) \cite{nvidia_jetson_orin_nano}.

Fällt das primäre Dual-GNSS-Signal in großer Höhe ($\SI{400}{\m}$ bis $\SI{4000}{\m}$ über Grund) aus, schaltet der Autopilot primär in den Modus der Koppelnavigation 
(\textit{Dead Reckoning}). Da die on-board integrierten, dreifach redundanten IMUs (Typ: \textit{InvenSense ICM-45686}) aufgrund mathematischer Integrationsfehler 
einer zeitlichen Positionsdrift von ca. $1\,\%$ bis $2\,\%$ der zurückgelegten Wegstrecke unterliegen, wird eine sensorische Stützung über den Companion Computer realisiert:

\begin{itemize}
    \item \textbf{Optischer Kartenabgleich (Map Matching / VIO):} Der \textit{NVIDIA Jetson Orin Nano} greift parallel zur KI-Inferenz auf den digitalen 4K-Videostrom 
                                                                der primären EO/IR-Nutzlast zu. Mittels Algorithmen zur \textit{Visual Inertial Odometry} (VIO) und lokal 
                                                                hinterlegten, georeferenzierten Satellitenbildern gleicht das System markante topografische Merkmale 
                                                                (z.~B. Bergkämme, Flussläufe oder Waldgrenzen) mit der Karte ab. Die hieraus ermittelten Positionskorrekturen 
                                                                werden in Echtzeit über die Ethernet-Infrastruktur an die FMU übermittelt, wodurch die Langzeitdrift der 
                                                                IMUs vollständig kompensiert wird. Die absolute Orientierung im Raum wird dabei kontinuierlich über das 
                                                                digitale Magnetometer (\textit{IST8310}) und die barometrische Höhe (\textit{BMP388}) gestützt.
\end{itemize}

Für den Übergang in den Schwebeflug sowie die anschließende autonome Vertikallandung unter GNSS-Verlust greift das UAS ab einer mittleren und bodennahen Flughöhe auf 
dedizierte Unterbodensensoren zurück:

\begin{itemize}
    \item \textbf{Eindimensionaler LiDAR-Entfernungsmesser (LightWare SF11/C):} Ab einer Höhe von \SI{120}{\m} über Grund erfasst der transienten-resistente 
                                                                                Infrarot-Laser-Entfernungsmesser präzise die wahre Höhe über Grund 
                                                                                (\textit{Altitude Above Ground Level} -- AGL) \cite{lightware_sf11c}. Der Sensor ist über 
                                                                                eine serielle UART-Schnittstelle direkt an das Pixhawk-Baseboard angebunden und liefert 
                                                                                mit einer Update-Rate von $\SI{20}{\Hz}$ driftfreie Höhendaten, die unbeeinflusst von 
                                                                                alpinen Luftdruckschwankungen sind.
    \item \textbf{Optischer Flusssensor (ARK Flow Modul):} Unterschreitet das Luftfahrzeug im Sinkflug eine Höhe von \SI{25}{\m}, klinkt sich das auf dem 
                                                           \textit{PMW3901}-Chipsatz basierende \textit{ARK Flow}-Modul ein \cite{ark_flow}. Der Sensor erfasst die 
                                                           horizontale Verschiebung von Bodenstrukturen mit hoher Frequenz. Die Datenübertragung erfolgt störsicher 
                                                           über das \textit{DroneCAN}-Protokoll via CAN-Bus. Dies ermöglicht es dem Extended Kalman Filter (EKF) des Autopiloten, 
                                                           jegliche windbedingte Horizontal-Drift im Schwebeflug ohne Satellitenverbindung vollständig zu eliminieren.
\end{itemize}

Das Zusammenspiel dieser Sensorkette sichert im SAIL-IV-Szenario ein kontrolliertes Abbremsen aus dem Vorwärtsflug, das Einleiten eines stabilen Schwebeflugs sowie eine 
zentimetergenaue, autonome Vertikallandung am Notlandeplatz.

\begin{table}[h!]
    \centering
    \caption{Sensoreinsatz und Redundanzstufen bei GNSS-Ausfall nach Flughöhe}
    \vspace{0.2cm}
    \begin{tabular}{p{3.2cm} p{4.2cm} p{4.6cm}}
        \toprule
        \textbf{Flughöhenbereich} & \textbf{Aktive Sensorik (Hardware)} & \textbf{Navigationsmethode / Failsafe} \\
        \midrule
        $\SI{400}{\m} - \SI{4000}{\m}$ & Redundante IMUs, Magnetometer, Baro, Suchkamera & IMU-Koppelnavigation gestützt durch optischen KI-Kartenabgleich (\textit{Jetson Orin Nano}) \\
        $\SI{25}{\m} - \SI{120}{\m}$ & Redundante IMUs, LiDAR (\textit{LightWare SF11/C}) & Höhengestützte IMU-Navigation zur Einleitung des autonomen Sinkflugs \\
        $\SI{0}{\m} - \SI{25}{\m}$ & LiDAR, Optical Flow (\textit{ARK Flow}) & \textit{Optical-Flow-Position-Hold} zur vollständigen Drifteliminierung bei der VTOL-Landung \\
        \bottomrule
    \end{tabular}
    \label{tab:navigation_failsafe_stages}
\end{table}

\subsubsection{Digital Airspeed Sensor (Pitot-Static Tube)}
\label{sec:airspeed_sensor}

Für den sicheren Starrflügler-Reiseflug sowie die kritischen Phasen des Transitionsflugs verfügt das UAS über ein digitales Differenzdrucksensorsystem zur Ermittlung der Eigengeschwindigkeit gegenüber der Luft (\textit{Indicated Airspeed} -- IAS). 
Das System nutzt das \textit{PT60 Pitot-Statik-Rohr} in Kombination mit dem industriellen \textit{Holybro High Precision DroneCAN Airspeed Sensor} \cite{holybro_dlvr_airspeed}. Mit einer Gesamtmasse von lediglich $38{,}3\,\text{g}$ (inkl. Verschlauchung) 
wird das Prandtl-Rohr an der äußersten Spitze der Rumpfnase positioniert, um Messfehler durch den \textit{Propwash} der Rotoren im Schwebeflug auszuschließen.

Die sensorische Erfassung und dezentrale Verarbeitung basiert auf folgenden Spezifikationen:
\begin{itemize}
    \item \textbf{DLVR-L10D Sensor mit CoBeam-Technologie:} Der Differenzdruckwandler nutzt die \textit{CoBeam\texttrademark}-Architektur von AllSensors. Diese garantiert eine Messgenauigkeit von besser als $1\,\%$ ($<1{,}0\,\%$ Fehler) über das 
                                                            gesamte Temperaturspektrum von $\SI{-20}{\celsius}$ bis $\SI{+85}{\celsius}$ und eliminiert thermische Nullpunktdriffts im alpinen Einsatz vollständig. Der Messbereich reicht bis \SI{2500}{\Pa}, 
                                                            was Fluggeschwindigkeiten bis $\SI{226.8}{\km/h}$ abbildet.
    \item \textbf{Dezentraler DroneCAN-Knoten:} Das Modul verfügt über einen dedizierten \textit{STM32G473CE} Mikroprozessor ($\SI{170}{\MHz}$, $512\,\text{KB}$ Flash) und führt die \textit{AP-Periph DroneCAN}-Firmware aus. 
                                                Die Druckwerte werden direkt auf Chipebene gefiltert und digital über das störsichere \textit{DroneCAN}-Protokoll (CAN-Bus) an die FMU übertragen. Dies schließt elektromagnetische Interferenzen (EMV) 
                                                durch benachbarte Hochstrom-Antriebskabel systembedingt aus.
\end{itemize}

Der Extended Kalman Filter (EKF) des Autopiloten fusioniert die hochfrequenten Airspeed-Daten permanent mit dem dualen GNSS-System zur Echtzeitberechnung des globalen Windvektors. Die Flugregelung nutzt diese Datenbasis als primäres 
Schaltkriterium im \textit{Conversion Mode}: Der Übergang in den Vorwärtsflug wird erst freigegeben, wenn der Sensor das stabile Erreichen der minimalen aerodynamischen Fluggeschwindigkeit (\textit{Stall Speed} plus Sicherheitsmarge) 
an den Regelkreis zurückmeldet.

\subsubsection{Flight Termination System (FTS) -- SAIL-IV Emergency Containment}
\label{sec:flight_termination_system}

Zur Erfüllung der behördlichen Auflagen hinsichtlich des Schutzes von Drittparteien am Boden (\textit{Enhanced Containment} gemäß EASA MOC 2511) verfügt das UAS über ein vom Autopiloten vollständig 
isoliertes Flugabbruchsystem vom Typ \textit{Dronavia Zephyr CC ESC MF} \cite{dronavia_zephyr_fts}. Während das im PX4-Autopiloten konfigurierte Software-Geo-Caging als erste Barriere dient und bei Erreichen eines 
harten Radius von $\SI{9.5}{\km}$ eine automatische Umkehr einleitet, fungiert das FTS als unüberbrückbare Hardware-Letztinstanz.

Das $\SI{9}{\g}$ leichte Modul wird im Rumpfzentrum in Reihe zwischen den PWM-Ausgängen des Pixhawks und den Eingängen der drei Haupt-ESCs geschaltet. Bei einer manuellen Aktivierung über den autarken LoRa-Bodensender 
(\textit{Klick-Fernauslöser}) fängt das FTS die Signale innerhalb von $\SI{0.27}{\s}$ ab und zwingt die Triebwerke in einen harten Not-Aus (\textit{Emergency Motor Cut-Off}). Zeitgleich wird über einen integrierten PWM-Ausgang 
das ballistische Fallschirmrettungssystem (PRS) gezündet. Die Kommunikation erfolgt verschlüsselt im $\SI{869}{\MHz}$-Band. Um im alpinen Gelände Signalabschattungen zu minimieren, wird die Bodensender-Antenne auf einem kompakten 
Stativ ($1{,}5\,\text{m}$) im Feld aufgebaut, was eine stabile Funkreichweite von bis zu $\SI{10}{\km}$ garantiert.

\subsubsection{Geometrically Isolated Backup Magnetometer Architecture}
\label{sec:isolated_magnetometer}

Aufgrund der hohen elektrischen Ströme ($\ge \SI{120}{\A}$) im zentralen Rumpfbereich während der VTOL-Startphase kommt es zu transienten elektromagnetischen Eigeninterferenzen. Diese verfälschen das im primären GNSS-Gehäuse (Holybro M10) integrierte Magnetometer temporär vollständig (\textit{Magnetic Interference}). 

Zur Gewährleistung einer ungestörten Kursreferenz wird ein sekundäres Hochpräzisions-Magnetometer vom Typ \textit{Holybro RM3100} eingesetzt \cite{holybro_RM3100}. Im Gegensatz zu klassischen AMR-Sensoren basiert das RM3100 auf einer 
temperaturstabilen, magneto-induktiven Messsensorik mit inhärenter Rauschfilterung. Der Sensor wird räumlich isoliert im äußeren Segment einer Tragfläche oder innerhalb des vertikalen Heckleitwerks verbaut, weit abseits aller Leistungskabel. 
Die Busanbindung erfolgt digital und störsicher via CAN-Bus über das \textit{DroneCAN}-Protokoll, wodurch Signaldämpfungen im Vergleich zu klassischen I2C-Verbindungen ausgeschlossen sind. Der Extended Kalman Filter (EKF) des Autopiloten gleicht beide 
Kompassquellen permanent ab und schaltet bei magnetischen Anomalien im Rumpfzentrum nahtlos und transientenfrei auf das isolierte Backup-Magnetometer um.

\subsubsection{Ground-to-Air Communication Infrastructure (Datalink Methodology)}
\label{sec:ground_to_air_communication}

Die Realisierung von BVLOS-Missionsradien innerhalb der behördlichen SAIL-IV-Zulassung erfordert eine redundante, hochbandbreitige Kommunikationsinfrastruktur zur simultanen Übertragung von Command-and-Control-Daten (C2), 
telemetrischen MAVLink-Zustandsvektoren und hochauflösenden Videostreams für die on-board Bildverarbeitung. Zur lückenlosen Abdeckung des theoretischen Einsatzszenarios sowie der praktischen Erprobungsphase des Prototragsystems wird 
ein zweistufiges Hardware-Konzept verfolgt, welches die wirtschaftlichen Limits eines schulischen Budgets mit den operativen Anforderungen harmonisiert.

\begin{enumerate}
    \item \textbf{Primäres Prototypen-Setup (Taktischer Feld-Link):} Für den operationalen Prototypen wird das digitale High-Definition-Übertragungssystem \textit{CubePilot Herelink v1.1} (Gebrauchtgerät-Spezifikation, Artikelnummer: MBS-DIP-12) 
    integriert \cite{mybotshop_herelink}. Das System operiert im \SI{2.4}{\GHz}-ISM-Band und kombiniert die Funktionen einer Fernsteuerung, eines Telemetrielinks und eines Videoübertragungssystems in einer kompakten Baugruppe. 
    Die Luftfahrzeugeinheit (\textit{Herelink Air Unit}) speist das H.264/H.265-komprimierte Full-HD-Videosignal ($1080\text{p}$ @ $60\,\text{fps}$) der EO/IR-Nutzlast über die hardwareseitig integrierte Ethernet-RJ45-Schnittstelle 
    direkt in den Datenstrom ein. 
    
    Die Signallaufzeit (\textit{End-to-End Glass-to-Glass Latenz}) beträgt herstellerseitig optimierte $\approx \SI{110}{\ms}$. Bei einer geringen Leistungsaufnahme von $< \SI{4}{\W}$ an Bord wird im direkten Sichtlinienbetrieb 
    (\textit{Line-of-Sight} -- LOS) eine stabile Funkreichweite von bis zu $\SI{20}{\km}$ erzielt. Die \textit{Herelink Ground Station} fungiert als primäre Steuereinheit für den Piloten und stellt über ein integriertes WLAN-Modul 
    einen lokalen Netzwerk-Hotspot zur Verfügung. Über diesen Hotspot greift der bodengestützte Missions-Laptop drahtlos auf den RTSP-Videostream für die Jetson-KI-Evaluierung sowie das volle \textit{QGroundControl}-Telemetriespektrum zu. 
    Zur Maximierung des Antennengewinns wird die Bodeneinheit stationär an einem kompakten Feldstativ ($1{,}5\,\text{m}$) mit omnidirektionalen und direktionalen Patch-Antennen betrieben.
    
    \item \textbf{Theoretisches Upgrade-Szenario (Militärischer Long-Range Link):} Da das primäre Herelink-System physikalisch auf eine maximale Reichweite von \SI{20}{\km} limitiert ist, erfordern Einsätze über diese Distanz hinaus ein technologisches Upgrade. 
    Für dieses erweiterte Anforderungsprofil wird theoretisch der Übergang zu industriellen IP-Mesh-Funkmodemen (z.~B. \textit{Microhard pDDL}-Serie) im \SI{2.4}{\GHz}-Bereich in Kombination mit einem autarken, schmalbandigen \textit{ExpressLRS (ELRS)} 
    LoRa-Funksystem im stark durchdringenden \SI{868}{\MHz}-UHF-Band definiert. Zur Kompensation der Freiraumdämpfung wird bodenseitig ein automatischer Antennentracker (\textit{Automatic Antenna Tracker} -- AAT) mit einer hochgewinnbringenden 
    Richtfunkantenne eingesetzt, welcher sich dynamisch nach den GNSS-Echtzeitkoordinaten des VTOLs ausrichtet.
\end{enumerate}

\begin{table}[h!]
    \centering
    \caption{Frequenzverteilung und funktionale Matrix des Kommunikations-Stacks}
    \vspace{0.2cm}
    \begin{tabular}{p{3.5cm} p{2.2cm} p{2.8cm} p{3.5cm}}
        \toprule
        \textbf{Subsystem} & \textbf{Frequenzband} & \textbf{Schnittstelle / Protokoll} & \textbf{Primäre Funktion (Einsatz)} \\
        \midrule
        CubePilot Herelink  & \SI{2.4}{\GHz} & COFDM / RTSP \& MAVLink & Video, Telemetrie \& C2 (Prototyp) \\
        Microhard pDDL-Serie & \SI{2.4}{\GHz} & IP-Mesh / TCP-UDP & High-Speed Datenträger ($>\SI{20}{\km}$ Option) \\
        ExpressLRS Link    & \SI{868}{\MHz} & CRSF / LoRa Modulation & Redundante C2-Steuerung (Option) \\
        \bottomrule
    \end{tabular}
    \label{tab:communication_matrix}
\end{table}

\newpage

\subsubsection{Regulatory Transponder Peripherals (Detect-and-Avoid \& Remote ID)}
\label{sec:regulatory_peripherals}

Zur Erlangung der behördlichen SAIL-IV-Fluggenehmigung im europäischen Luftraum müssen zwingend Systeme zur kooperativen Luftraumüberwachung sowie zur digitalen Identifikation integriert werden. Diese Komponenten minimieren das Risiko von 
Kollisionen im BVLOS-Sichtlinienbetrieb und stellen die gesetzliche Konformität sicher:

\begin{itemize}
    \item \textbf{ADS-B In Empfänger (uAvionix pingRX Pro):} Für die Umsetzung eines automatisierten \textit{Detect-and-Avoid}-Konzepts (DAA) wird ein Luftfahrt-zertifizierter Dual-Band-ADS-B-Empfänger vom Typ \textit{uAvionix pingRX Pro} 
    integriert \cite{uavionix_pingrx_pro}. Das System wiegt lediglich $5\,\text{g}$ bei minimalen Gehäuseabmessungen von $32 \times 31 \times 9\,\text{mm}$ und schützt seine Elektronik über ein robustes Aluminiumgehäuse. 
    Der Empfänger wertet simultan die Frequenzen $\SI{978}{\MHz}$ und $\SI{1090}{\MHz}$ aus, um Transpondersignale bemannter Luftfahrzeuge (z.~B. Rettungshubschrauber) in Echtzeit zu erfassen. Die Anbindung erfolgt über ein standardisiertes 
    MAVLink-Protokoll direkt an einen freien Telemetrie-UART-Port der FMU. Die empfangenen Positionsdaten fremder Luftfahrzeuge werden direkt in die Luftraumkarte der GCS eingespeist, um dem Piloten eine präzise Lageerkennung zu ermöglichen.
    
    \item \textbf{Direct Remote ID Modul (Holybro Remote ID Module):} Zur Erfüllung der gesetzlichen EU-Fernkennungsverordnung verfügt das UAS über ein dediziertes \textit{Holybro Remote ID Module} \cite{holybro_remote_id}. 
    Das Modul basiert auf einem \textit{Espressif ESP32-C3}-Funkchip und strahlt die Betreiber-ID sowie die Live-Telemetriedaten (Position, Höhe, Kurs) mit einer Sendeleistung von $+20\,\text{dBm}$ ($100\,\text{mW}$) via $\SI{2.4}{\GHz}$ 
    WLAN und Bluetooth aus, wodurch im Freifeld Detektionsreichweiten von $>\SI{5}{\km}$ erzielt werden. Das System führt die Open-Source-Firmware \textit{ArduRemoteID} aus und ist softwareseitig nativ im PX4-Autopiloten verankert. 
    Die physische Einbindung erfolgt im Avionikdeck, wobei die Spannungsversorgung ($\SI{5}{\V}$) und der Datenaustausch verlustfrei über den digitalen CAN-Bus mittels 
    des \textit{DroneCAN}-Protokolls realisiert werden.
\end{itemize}

\subsection{On-Board Communication Busses (Bus-Topologien)}
\subsubsection{Deterministic Control via CAN-Bus (UAVCAN / Cyphal)}
\subsubsection{High-Bandwidth Inter-Processor Bridge (Ethernet)}

\subsection{Power Distribution Architecture \texorpdfstring{\&}{\&} Safety Isolation}
\subsubsection{Redundant Avionics Power Supply (BECs)}
\subsubsection{Galvanic Isolation of the Propulsion System}

\subsection{Electromagnetic Compatibility (EMV) \texorpdfstring{\&}{\&} Shielding Concept}
\subsubsection{Component Placement and Spatial Separation}
\subsubsection{Physical Shielding and Grounding}


\newpage
\section{Payload Weight Statement (Massen- und Schwerpunktbilanz)}
\label{sec:weight_balance}

Die exakte Erfassung aller Einzelmassen innerhalb des Avionik-, Sensorik- und Nutzlast-Stacks ist eine fundamentale Voraussetzung für die Bestimmung des Gesamtschwerpunkts (\textit{Center of Gravity} -- CoG) sowie für die aerodynamische Trimmung des Tri-Tilt-Rotor-UAS. Da das System ein modulares Payload-Konzept für die Suchkamera verfolgt, bildet die Massenbilanz in Tabelle~\ref{tab:avionics_mass_balance} das maximale Gewichtsszenario (unter Verwendung des schwersten \textit{Foxtech Seeker AI}-Gimbals inklusive mechanischer Dämpfungsadapter) ab.

\begin{table}[h!]
    \centering
    \caption{Vollständige Massenbilanz des Avionik-, Sensorik- und Peripherie-Stacks}
    \vspace{0.2cm}
    \begin{tabular}{p{5.8cm} p{4.7cm} p{2.5cm}}
        \toprule
        \textbf{Baugruppe / Komponente} & \textbf{Spezifikation / Hardware-Typ} & \textbf{Masse (g)} \\
        \midrule
        \textbf{Flugsteuerung (Core)} & & \\
        ~~Pixhawk 6X Autopilot & Flight Management Unit (FMU, STM32H753) & $\SI{120}{\g}$ \\
        ~~Holybro Standard Baseboard & v2B-Trägerplatine mit Ethernet-Logik & $\SI{43}{\g}$ \\
        \midrule
        \textbf{Missionscomputer} & & \\
        ~~NVIDIA Jetson Orin Nano & System-on-Module (SoM) & $\SI{30}{\g}$ \\
        ~~Reference Carrier Board & Inkl. Aluminium-Kühlkörper \& Aktivlüfter & $\SI{125}{\g}$ \\
        \midrule
        \textbf{Satellitennavigation} & & \\
        ~~Holybro M10 GPS (Primär) & u-blox M10 Chipsatz, inkl. IST8310 Kompass & $\SI{42}{\g}$ \\
        ~~Holybro H-RTK F9P Rover (Sekundär) & u-blox ZED-F9P Multiband-Empfänger & $\SI{36}{\g}$ \\
        \midrule
        \textbf{Integrierte Autopilot-Sensoren} & & \\
        ~~3x InvenSense ICM-45686 & Dreifach redundante IMUs (auf FMU-Board) & im Core inkl. \\
        ~~Diversitäre Barometer-Bestückung & 1x ICP20100 \& 1x BMP388 (auf FMU-Board) & im Core inkl. \\
        \midrule
        \textbf{Missions- \& Failsafe-Sensorik} & & \\
        ~~Holybro High Precision DroneCAN & DLVR Differenzdruck-Wandlerplatine & $\SI{16}{\g}$ \\
        ~~Prandtl-Staurohr (PT60) & Pitot-Statik-Sonde inkl. Verschlauchung & $\SI{22.3}{\g}$ \\
        ~~LightWare SF11/C & Eindimensionaler LiDAR-Entfernungsmesser & $\SI{35}{\g}$ \\
        ~~ARK Flow Modul & PMW3901 Optical Flusssensor via DroneCAN & $\SI{12}{\g}$ \\
        ~~Holybro RM3100 & Geometrisch isolierter Backup-Kompass & $\SI{18}{\g}$ \\
        \midrule
        \textbf{Kommunikation \& Regulierungen} & & \\
        ~~CubePilot Herelink Air Unit v1.1 & Primärer C2- \& HD-Videodatenlink & $\SI{95}{\g}$ \\
        ~~2x Herelink Taoglas Antennen & Luftfahrzeug-Funkantennen ($\SI{2.4}{\GHz}$) & $\SI{30}{\g}$ \\
        ~~Herelink Air Unit Ethernet-Kabel & Spezialadapter auf RJ45-Netzwerkbuchse & $\SI{12}{\g}$ \\
        ~~Dronavia Zephyr CC ESC MF & Dediziertes Hardware-FTS Modul & $\SI{9}{\g}$ \\
        ~~uAvionix pingRX Pro & Dual-Band ADS-B-In-Empfänger & $\SI{5}{\g}$ \\
        ~~Holybro Remote ID Module & EU-Fernkennungs-Platine (ohne Gehäuse) & $\SI{4}{\g}$ \\
        \midrule
        \textbf{Nutzlast \& Verkabelung} & & \\
        ~~Foxtech Seeker EO \& Ranging AI & 3-Achsen-Gimbal-Kamera (Max-Szenario) & $\SI{800}{\g}$ \\
        ~~Payload Interface \& Dämpfung & Mechanische CFK-Aufhängung \& Puffer & $\SI{65}{\g}$ \\
        ~~Avionik-Gesamtverkabelung & JST-GH, Cat.6-Ethernet, Coaxial \& CAN & $\SI{150}{\g}$ \\
        \midrule
        \textbf{Gesamtgewicht Avionik/Payload} & \textbf{Reine Elektronik- und Nutzlastmasse} & \textbf{\SI{1648.3}{\g} (\SI{1.648}{\kg})} \\
        \bottomrule
    \end{tabular}
    \label{tab:avionics_mass_balance}
\end{table}


\section{Flight Control Laws (C-Laws for Conversion Mode)}
\section{Sensor Fusion Matrix (IMU, GNSS, Air-Data)}
\chapter{Electrical \& Energy Storage Systems}
\label{ch:electrical_energy}

Dieses Kapitel behandelt die elektrotechnische Architektur, die Energiespeicherkomponenten sowie die ausfallsicheren Versorgungsnetzwerke des VTOL-Systems. Um die Betriebssicherheit des Gesamtsystems unter extremen atmosphärischen Bedingungen zu gewährleisten, wird eine strikte physische und galvanische Systemtrennung zwischen den Hochstrom-Antriebskreisen und den sensitiven Niedervolt-Avioniksystemen realisiert.

\noindent Der Fokus der Auslegung liegt auf der Implementierung eines intelligenten Batterie-Management-Systems (BMS) zur permanenten Einzelzellenüberwachung sowie auf der Definition konstruktiver thermischer Schutzmaßnahmen für den alpinen Worst-Case-Einsatz. Ergänzt wird die Energiearchitektur durch ein dual-redundantes Bus-Backup-System, welches bei einem totalen Ausfall der Hauptstromversorgung die unterbrechungsfreie Speisung aller flugsicherheitskritischen Steuer- und Aktuatorsysteme autonom sicherstellt.

\section{Power Distribution Architecture (High- vs. Low-Volt)}
\section{Battery Management System (BMS) \& Thermal Management}
\section{Emergency Power (Bus-Backup)}

\chapter{Communication \& Datalink}
\label{ch:communication}
\section{Command \& Control (C2) Datalink}
\section{Payload Datalink}
\section{Lost-Link Procedures (Autonome RTH-Routinen)}

\chapter{Payload Integration}
\label{ch:payload}
\section{Payload Bay Interfaces (Mechanical \& Data)}
\section{Center of Gravity Impact}

\chapter{Safety, Redundancy \& Failure Modes (FMEA)}
\label{ch:safety}
\section{Functional Hazard Assessment (FHA)}
\section{Redundancy Concept (Fail-Safe vs. Fail-Operational)}
\section{Emergency Systems (Ballistischer Fallschirm)}

\chapter{Ground Control Station (GCS) \& HMI}
\label{ch:gcs_hmi}
\section{GCS Hardware \& Software Architecture}
\section{Pilot Interface (HMI \& Telemetry Visualization)}

\chapter{Manufacturing, Assembly \& Maintenance}
\label{ch:manufacturing}
\section{Design for Manufacturing (DFM)}
\section{Maintenance \& Component Lifespan (MTBF)}

\chapter{Verification, Validation \& Testing (V\&V)}
\label{ch:testing}
\section{Simulation \& Hardware-in-the-Loop (HITL)}
\section{Ground Testing (Wing-Bend, Propulsion Testbed)}
\section{Flight Test Program (Transition Corridor Validation)}


\end{document}